% answers_short.tex —— 1987 年·全国卷·理科 · 简短答案速查
% 由 tools/make_short_answers.py 自 answers.tex 的 \gksolhead/\gkanswer 提取，
% 并与 1987/1987_数学真题_全国理_答案要点.md 逐题比对后生成，勿手改（改源 MD 后重跑）。

\begin{gknotice}
\textbf{使用说明}：本节只给最终答案便于快速核对，完整推导见解析版。答案为编写者独立推导，\textbf{非官方标准答案}。
\end{gknotice}

\gkpart{题目}
\begin{itemize}[leftmargin=2.6em]
  \item[\textbf{1.}] D
  \item[\textbf{2.}] C
  \item[\textbf{3.}] B
  \item[\textbf{4.}] A
  \item[\textbf{5.}] B
  \item[\textbf{6.}] D
  \item[\textbf{7.}] B
  \item[\textbf{8.}] A
  \item[\textbf{9.}] $\dfrac{3\pi}{2}$
  \item[\textbf{10.}] $\lambda<-2$ 或 $\lambda>-1$
  \item[\textbf{11.}] $8$
  \item[\textbf{12.}] $2$
  \item[\textbf{13.}] $\left(\dfrac12,1\right)$
  \item[\textbf{14.}] $72$
  \item[\textbf{15.}] $3\pi$
  \item[\textbf{16.}] $\dfrac1{16}$
  \item[\textbf{17.}] $V=\dfrac16L^2h$
  \item[\textbf{18.}] $0<a<1$
  \item[\textbf{19.}] (1) $|z_1+A||z_2+A|=|A|^2$；(2) $\displaystyle\frac{z_1+A}{z_2+A}=\left|\frac{\overline{z_1}+\overline A}{\overline{z_2}+\overline A}\right|$
  \item[\textbf{20.}] (1) $\displaystyle (1+b)a_n=ba_{n-1}+\frac{b}{(1+b)^n}$；(2) $\displaystyle a_n=\frac{b(1+b+\cdots+b^{n-1})}{(1+b)^{n+1}}$；(3) $\displaystyle\lim_{n\to\infty}S_n=1\quad(0<b<1)$
  \item[\textbf{21.}] 到 $y$ 轴的最短距离为 $\dfrac54$；此时 $M$ 为 $\left(\dfrac54,\dfrac{\sqrt2}{2}\right)$ 或 $\left(\dfrac54,-\dfrac{\sqrt2}{2}\right)$。
  \item[\textbf{22.}] 固定源下标为 $n\to\infty$，但式中底数和指数为 $x$，按字面无法得到唯一的关于 $n$ 的极限；若按通常意图统一改为 $\displaystyle\lim_{n\to\infty}\left(1-\frac1{2n}\right)^n$，则值为 $e^{-1/2}$。
  \item[\textbf{23.}] $\displaystyle y'=\ln(1+x^2)+\frac{2x^2}{1+x^2}$
\end{itemize}
